\section*{Chapitre \ref{ch:b3:rna-regulation} --- ARN non codants et régulation post-transcriptionnelle}

\begin{solution}{exo:b3:rna-regulation:1}
miARN~: codé dans le génome sous forme de tige-boucle, \qty{22}{nt},
découpé par \omterm{def:b3:rna-regulation:mirna}{Drosha} puis par \omterm{def:b3:rna-regulation:rnai}{Dicer}, s'apparie imparfaitement par sa
graine et réprime la traduction et la stabilité de nombreux messagers.
\omterm{def:b3:rna-regulation:ncrna}{siARN}~: issu d'un long ARN double brin (viral, transgénique,
expérimental), \qty{21}{nt}, dépendant de \omterm{def:b3:rna-regulation:rnai}{Dicer}, s'apparie parfaitement
et dirige le clivage de la cible par \omterm{def:b3:rna-regulation:rnai}{Argonaute}. \omterm{def:b3:rna-regulation:ncrna}{piARN}~: issu de
transcrits simple brin d'\omterm{def:b3:rna-regulation:pirna}{amas de piARN}, \qtyrange{26}{31}{nt},
indépendant de \omterm{def:b3:rna-regulation:rnai}{Dicer}, chargé sur les protéines PIWI dans la lignée
germinale, clive les transcrits de transposons (ping-pong) et dirige la
répression chromatinienne de l'ADN des transposons.
\end{solution}

\begin{solution}{exo:b3:rna-regulation:2}
La Pol~II transcrit le pri-miARN~; \omterm{def:b3:rna-regulation:mirna}{Drosha} (avec DGCR8) découpe la
tige-boucle dans le noyau~; l'exportine-5 exporte le pré-miARN~; \omterm{def:b3:rna-regulation:rnai}{Dicer}
retire la boucle et laisse un duplex de \qty{22}{nt}~; un brin est
chargé dans \omterm{def:b3:rna-regulation:rnai}{Argonaute}~; la graine (nucléotides 2--8) s'apparie à un
site de la 3$'$ UTR~; \omterm{def:b3:rna-regulation:rnai}{Argonaute} recrute la machinerie de désadénylation
et de décoiffage et bloque l'initiation --- le messager est moins
traduit et se dégrade plus vite.
\end{solution}

\begin{solution}{exo:b3:rna-regulation:3}
L'ARN fabriqué in vitro par une polymérase de phage contenait des
contaminants double brin (la polymérase copie aussi l'autre brin de la
matrice et déborde des extrémités), de sorte que les préparations
«~sens~» portaient le vrai déclencheur. Fire et Mello ont purifié
chaque brin, montré que chacun seul faisait peu, puis les ont appariés
délibérément~: l'ARN double brin était au moins cent fois plus
puissant, quelques molécules par cellule suffisant.
\end{solution}

\begin{solution}{exo:b3:rna-regulation:4}
Privée de fer~: l'IRP se fixe sur les IRE. Ferritine~: l'initiation de
la traduction est bloquée, aucune protéine de stockage n'est produite
(contrôle de la traduction). Récepteur de la transferrine~: le messager
est protégé de sa nucléase et s'accumule, donc plus de récepteur est
produit et plus de fer importé (contrôle de la stabilité du messager).
\end{solution}

\begin{solution}{exo:b3:rna-regulation:5}
$\gamma = \ln 2/\qty{120}{min} = \qty{0.0058}{min^{-1}}$~; $m^{*} =
5/0.0058 \approx 870$ molécules. Dégradation doublée~: $m^{*} \approx 430$,
temps de demi-approche $\qty{120}{min}/2 = \qty{60}{min}$.
\end{solution}

\begin{solution}{exo:b3:rna-regulation:6}
Séquence 3$'$ UTR totale $2\times 10^{7}$ nt~; correspondances dues au hasard $2\times
10^{7}/4^{7} = 2\times 10^{7}/\num{16384} \approx \num{1200}$ sites.
Une vraie cible se reconnaît à la conservation du site entre espèces,
à un contexte favorable (voisinage riche en AU, site éloigné du codon
stop, appariement supplémentaire à l'extrémité 3$'$ du miARN), à la
coexpression du miARN et de la cible dans la même cellule, et à
l'expérience~: le messager baisse quand le miARN est présent et ne
baisse plus quand le site est muté.
\end{solution}

\begin{solution}{exo:b3:rna-regulation:7}
(a) XX, \emph{tra} nul~: SXL est produite mais pas TRA, \emph{dsx} prend
l'épissage mâle par défaut --- un mâle somatique (stérile). (b) XY avec
TRA constitutive~: TRA-2 est présente dans les deux sexes, donc DSX est
produite sous la forme femelle --- développement somatique femelle (une
«~pseudo-femelle~» stérile, la lignée germinale suivant d'autres
indices). (c) Pas de \emph{dsx}~: aucune des deux formes de DSX, donc
aucun des deux programmes de différenciation terminale n'est imposé ---
un intersexué aux caractères mixtes ou incomplets.
\end{solution}

\begin{solution}{exo:b3:rna-regulation:8}
À l'état stationnaire le messager, et donc la protéine, est
proportionnel à la demi-vie du messager~:
$\qty{180}{min}/\qty{15}{min} = 12$ fois plus de protéine. La protéine est un signal de croissance
normal censé être transitoire, une impulsion après chaque stimulus~;
produite douze fois plus et en continu, elle entraîne la prolifération
sans stimulus. C'est la dose et la chronologie, non la séquence, qui la
rendent oncogène.
\end{solution}

\begin{solution}{exo:b3:rna-regulation:9}
(a) Mère de laboratoire, père sauvage~: les œufs ne portent pas de
\omterm{def:b3:rna-regulation:ncrna}{piARN} contre les éléments P (la lignée de la mère ne les a jamais
rencontrés), les éléments P paternels transposent librement dans la
lignée germinale de la descendance --- stérile (dysgénésique). (b) Mère
sauvage~: ses œufs contiennent des \omterm{def:b3:rna-regulation:ncrna}{piARN} d'éléments P, qui répriment
les éléments chez la descendance --- fertile. L'asymétrie tient au
dépôt maternel des \omterm{def:b3:rna-regulation:ncrna}{piARN}, hérédité d'un ARN et non de gènes.
\end{solution}

\begin{solution}{exo:b3:rna-regulation:10}
La protéine intègre son messager sur sa propre durée de vie
$\tau_{p}$. Les messagers vivent $\tau_{m} = 1/(\gamma + \gamma' R)$~;
pendant $\tau_{p}$ la protéine voit donc environ $\tau_{p}/\tau_{m}$
durées de vie de messager indépendantes, chacune un événement aléatoire
de naissance et de mort. La fluctuation relative d'une moyenne sur $N$
événements indépendants décroît en $1/\sqrt{N}$, de sorte qu'un
$\tau_{m}$ plus court sous contrôle d'un \omterm{def:b3:rna-regulation:ncrna}{microARN} --- à messager moyen
égal, ce qui exige un $k$ proportionnellement plus élevé --- donne une
protéine plus lisse. La cellule paie le \omterm{def:b3:rna-regulation:ncrna}{microARN} et la transcription
supplémentaire pour acheter de la précision et de la vitesse~: un
niveau fixé par un équilibre entre synthèse rapide et dégradation
rapide est à la fois moins bruité et plus prompt à se réajuster que le
même niveau fixé par une synthèse lente et une dégradation lente.
\end{solution}

\begin{solution}{exo:b3:rna-regulation:11}
Ver~: l'ARN polymérase ARN-dépendante fabrique de nouveaux \omterm{def:b3:rna-regulation:ncrna}{siARN} à
partir du messager cible lui-même, de sorte que le signal est renouvelé
tant que la cible est transcrite, survit à la dilution par les
divisions et, grâce à un canal à ARN double brin entre cellules, gagne
d'autres tissus et la lignée germinale pour quelques générations.
Mammifère~: pas d'amplification, donc les \omterm{def:b3:rna-regulation:ncrna}{siARN} sont consommés et
dilués à chaque division, n'agissent que dans les cellules qui les ont
reçus, et s'estompent en quelques jours dans des cellules qui se
divisent. Conséquence~: un médicament à \omterm{def:b3:rna-regulation:ncrna}{siARN} doit être porté dans
chaque cellule cible, agit au mieux dans des cellules qui ne se
divisent pas (hépatocytes, d'où les médicaments hépatiques) et demande
des administrations répétées.
\end{solution}

\begin{solution}{exo:b3:rna-regulation:12}
(a) ARN fonctionnel~: détruire le transcrit après sa production (\omterm{def:b3:rna-regulation:ncrna}{siARN},
oligonucléotide antisens) donne un phénotype, et exprimer l'ARN depuis
un transgène placé ailleurs dans le génome le corrige. (b)
Transcription fonctionnelle~: insérer un signal de polyadénylation
précoce qui arrête la transcription (ou supprimer le promoteur) donne
le phénotype, alors que détruire l'ARN après coup ne le donne pas, et
un transgène ectopique ne corrige rien --- c'est l'acte de
transcription (ou l'élément d'ADN) qui compte, comme pour bien des ARN
associés aux amplificateurs. (c) Bruit~: aucune manipulation ne change
rien, l'ARN est à moins d'une copie par cellule, sa séquence et son
expression ne sont pas conservées.
\end{solution}

\begin{solution}{pb:b3:rna-regulation:1}
\textbf{1.} $\gamma = \ln 2/40 = \qty{0.0173}{min^{-1}}$~; $m^{*} =
3/0.0173 \approx 170$ messagers.
\textbf{2.} $\delta_{p} = \ln 2/180 = \qty{0.00385}{min^{-1}}$~; $P^{*} =
\beta m^{*}/\delta_{p} = 2\times 173/0.00385 \approx 9.0\times 10^{4}$
protéines.
\textbf{3.} Dégradation $4\gamma = \qty{0.069}{min^{-1}}$~; $m^{*} = 3/0.069
\approx 43$~; constante de temps $1/0.069 = \qty{14}{min}$.
\textbf{4.} $\beta = 1$~: $P^{*} = 43/0.00385 \approx 1.1\times 10^{4}$~;
répression d'un facteur $8$ ($4$ par la dégradation, $2$ par la
traduction).
\textbf{5.} Demi-approche du messager $\ln 2/(4\gamma) = \qty{10}{min}$.
La protéine décroît ensuite avec sa propre demi-vie, \qty{3}{h}~: c'est
la dégradation de la protéine qui limite le basculement.
\textbf{6.} Oui~: le messager ne tombe qu'à $1/4$ (encore aisément
détecté) tandis que la protéine tombe à $1/8$ et continue de baisser à
mesure que l'ancienne protéine se renouvelle~; l'observation
ancienne d'une protéine disparue «~alors que l'ARNm demeurait~»
reflétait surtout la moitié traductionnelle de la répression.
\textbf{7.} $10^{6}/250 = 4000$ molécules par œuf.
\textbf{8.} $4000/550 \approx 7$ par cellule à l'éclosion. $4000/2^{n} < 1$
pour $n \ge 12$ divisions depuis le zygote --- environ trois divisions
après l'éclosion.
\textbf{9.} La polymérase recopie le messager cible en un nouvel ARN
double brin, découpé en \omterm{def:b3:rna-regulation:ncrna}{siARN} secondaires~: le déclencheur est
régénéré depuis la cible elle-même, donc il survit à la dilution et se
propage (un canal fait passer l'ARN double brin entre cellules et dans
la lignée germinale). Il s'estompe parce que l'amplification exige à la
fois la cible et un déclencheur primaire~; une fois l'ARN injecté
disparu, le stock secondaire se dilue de génération en génération et la
lignée germinale se remet à zéro.
\textbf{10.} $5000/2^{n} < 100$~: $2^{n} > 50$, $n = 6$ jours ($78$
copies restantes).
\textbf{11.} Des cellules qui ne se divisent pas ne diluent pas le \omterm{def:b3:rna-regulation:rnai}{RISC}~;
la limite est la durée de vie chimique du \omterm{def:b3:rna-regulation:ncrna}{siARN} (dégradation par les
nucléases, ralentie par la modification chimique des brins) et la lente
fuite du dépôt endosomal qui alimente le cytoplasme.
\textbf{12.} $2\times 10^{7}/\num{16384} \approx \num{1200}$
correspondances dues au hasard. C'est la génétique qui a tranché~: la
perte de \emph{lin-14} donne le phénotype inverse de la perte de
\emph{lin-4}, les allèles de gain de fonction de \emph{lin-14} étaient
des délétions portant exactement sur les sites de la 3$'$ UTR, et les
sept sites étaient conservés.
\textbf{13.} Niveau $\propto$ demi-vie~: hausse d'un facteur $50/10 = 5$.
\textbf{14.} Chute d'un facteur $1/0.05 = 20$ de la synthèse de ferritine.
\textbf{15.} Le rapport importation/stockage monte d'un facteur
$5\times 20 = 100$ entre l'état riche et l'état pauvre en fer.
\textbf{16.} Le centre ne s'assemble que si du fer cytosolique est
disponible, de sorte que la conformation de la protéine est une lecture
directe de la concentration en fer, sans récepteur, sans hormone ni
transcription --- comme un riborégulateur, un métabolite perçu par la
molécule qu'il commande.
\textbf{17.} La ferritine est traduite quel que soit le fer~: ferritine
sérique très élevée, fer sérique et réserves normaux (pas de
surcharge). Dans le cristallin, qui ne se renouvelle pas, la ferritine
s'accumule des années et cristallise --- une cataracte précoce.
\textbf{18.} La tige-boucle n'est qu'un site d'accueil~; l'effet est
tout ce que la protéine fixée gêne physiquement --- les ribosomes qui
balaient l'extrémité 5$'$, une nucléase à l'extrémité 3$'$.
\textbf{19.} Deux chimies couvrent deux conditions~: le centre
fer--soufre d'IRP1 est aussi détruit par les oxydants et par un
dioxygène rare, de sorte qu'IRP1 répond à l'état d'oxydoréduction,
tandis que la dégradation d'IRP2 dépendante du fer répond au fer
lui-même sur toute la gamme physiologique de dioxygène~; la redondance
rend le régulon robuste, et les deux capteurs rapportent des choses
différentes.
\textbf{20.} $12\times 48\times 33\times 2 = \num{38016}$. $1 - (1 -
50/\num{38016})^{50} = 1 - 0.99868^{50} \approx 1 - e^{-0.066} =
0.064$~: environ \qty{6}{\%}.
\textbf{21.} Des isoformes identiques se lient de façon homophile et se
repoussent~: les branches d'un même neurone s'évitent et s'étalent,
tandis que les branches de neurones différents, qui ne partagent
presque aucune isoforme, peuvent se croiser --- une identité privée par
cellule.
\textbf{22.} $2^{n}$~; $\num{1024}$ pour $n = 10$. Bien moins en
réalité~: l'inclusion est décidée par des facteurs propres au type
cellulaire et coordonnée entre exons, beaucoup de combinaisons
décalent le cadre de lecture et sont détruites par la dégradation des
ARNm non-sens, et les exons ne sont pas indépendants.
\textbf{23.} \emph{Sxl} commande aussi la \omterm{def:b3:chromatin-epigenetics:x-inactivation}{compensation de dose}~: chez la
femelle, SXL bloque la traduction de \emph{msl-2}, de sorte que les
deux chromosomes~X ne sont pas hypertranscrits~; sans SXL un embryon XX
assemble le complexe MSL sur ses deux chromosomes~X et double leur
production --- surdose létale de tous les gènes liés à l'X, avant même
que le sexe soit en jeu.
\textbf{24.} Le promoteur précoce ne fonctionne que tant que le signal
X:A est présent~; la SXL qu'il produit impose l'épissage productif du
transcrit issu du promoteur de maintien, actif dans les deux sexes, de
sorte qu'une fois SXL existante elle continue de se produire elle-même
et le signal n'est plus nécessaire --- la même logique auto-entretenue
du lecteur qui recrute l'écrivain que pour les marques de \omterm{def:b3:chromatin-epigenetics:nucleosome}{chromatine},
ici avec une protéine et un épissage.
\textbf{25.} Protéine LIN-14 réprimée d'un facteur $8$~; demi-basculement
du messager \qty{10}{min}~; un ARN non amplifié tombe sous une molécule
par cellule après $12$ divisions.
\end{solution}
